\section{Introduction}
\label{sec:introduction}

Large language models (LLMs) have become increasingly capable at code-related tasks. From automated program repair to repository-level code generation, LLM-powered code agents are reshaping modern software development at an unprecedented pace \citep{yang2024swe, jimenez2024swebench}. As these systems take on greater responsibility in safety-critical pipelines --- generating code that controls infrastructure, handles financial transactions, and manages medical records --- a fundamental question arises: do we actually understand how LLMs process code internally? High accuracy on downstream benchmarks provides a measure of \textit{what} these models can do, but tells us remarkably little about \textit{how} they do it. When a model correctly evaluates a loop with nested conditionals, is it genuinely tracking variable states and reasoning about execution paths, or is it exploiting surface-level shortcuts invisible to end-task evaluation? This black-box status limits our ability to diagnose failures, predict unreliable behaviors, and build trustworthy code agents.

To move beyond behavioral benchmarking, we need a structured framework for dissecting code comprehension. We draw on a foundational decomposition from programming language theory and compiler design: the separation of \textbf{control flow analysis} --- reasoning about execution paths, branch selection, and loop termination --- and \textbf{data flow analysis} --- tracking how values are produced, propagated, and consumed through variables \citep{aho2006compilers, allen1970control, kildall1973unified, nielson1999principles}. This decomposition, which has organized program analysis for over five decades, provides a principled axis along which to probe whether LLMs develop distinct internal mechanisms for different facets of code reasoning. We design five task families along this axis --- spanning pure variable tracking, arithmetic computation, conditional branching, loop execution, and short-circuit evaluation --- each parameterized for systematic difficulty control, with answers constrained to single-digit outputs to enable clean layer-wise analysis.

Most existing interpretability methods ask \textit{what information} is encoded in a model's hidden states --- probing for linguistic properties \citep{belinkov2017neural}, decomposing features via sparse autoencoders \citep{cunningham2023sparse}, or projecting activations into vocabulary space through logit lens variants \citep{nostalgebraist2020logit, belrose2023eliciting}. These approaches characterize \textit{what the model knows} at a given layer. We ask a different, more direct question: \textbf{at each layer, has the model already solved the problem?} This shift from representational content to behavioral trajectory leads us to a dual-bound diagnostic framework. We combine linear probing --- which tests whether the correct answer is \textit{linearly separable} in the hidden state, serving as a necessary condition for the model possessing the information --- with activation verbalization \citep{ghandeharioun2024patchscopes}, which tests whether the model can \textit{decode} the correct answer from the hidden state on its own, serving as a sufficient condition for information being fully formed. The gap between these two bounds --- information that exists but is not yet decodable --- becomes our primary analytical signal.

What if the hardest part of code reasoning isn't finding the answer --- but holding onto it? Applying our framework across five model scales (0.5B to 14B parameters), we discover that LLMs undergo a characteristic \textbf{brewing} phase during which internal representations transition from chaotic, high-entropy states to convergent, low-entropy distributions. This brewing phase terminates in one of four \textbf{resolution outcomes}: (1)~\textbf{Resolved} --- the model converges to the correct answer and maintains it through the remaining layers; (2)~\textbf{Overprocessed} --- the correct answer emerges at intermediate layers but is destroyed by continued computation, a form of self-inflicted failure; (3)~\textbf{Misresolved} --- the model converges confidently to a wrong answer, the most dangerous failure mode; and (4)~\textbf{Unresolved} --- the model never converges, running out of layers before completing its computation. Crucially, these resolution outcomes exhibit distinct distributional signatures (entropy profiles, inter-layer KL divergence, probing--verbalization agreement) that can be detected \textit{without access to ground-truth labels}, achieving approximately 95\% classification accuracy on our benchmark. Different code reasoning primitives trigger different distributions over these outcomes --- revealing that ``code understanding'' is not a monolithic capability but a collection of sub-processes with distinct internal dynamics.

In this paper, we develop a layer-wise diagnostic framework that traces an LLM's internal computation to determine whether it has already solved a code reasoning problem, is still brewing toward an answer, or has silently committed to a wrong one. Our contributions are threefold:
\begin{itemize}
    \item We propose a \textbf{dual-bound diagnostic framework} that combines probing (information separability) and activation verbalization (information decodability) as complementary necessary/sufficient conditions, using their gap as a first-order analytical signal for characterizing layer-wise reasoning trajectories.
    \item We construct an \textbf{interpretability-oriented code reasoning benchmark} of five task families parameterized along the control-flow / data-flow axis, co-designed with the diagnostic framework to ensure clean analytical signals.
    \item We report \textbf{systematic findings} on the internal reasoning dynamics of LLMs: four resolution outcomes that emerge from a universal brewing phase, each detectable without ground truth; task-dependent behavioral differences across code reasoning primitives; and causal validation through re-injection experiments showing that 80--90\% of unresolved cases can be rescued by extending computation, confirming that these models often possess reasoning capacity that their fixed depth fails to fully utilize.
\end{itemize}
Our findings have direct implications for building more reliable code agents and for the broader understanding of how LLMs manage --- and mismanage --- structured reasoning.
